\documentclass[11pt]{article}

\usepackage{amsmath,amssymb}
\usepackage{booktabs}
\usepackage{geometry}
\usepackage{hyperref}
\usepackage[numbers,sort&compress]{natbib}

\geometry{margin=1in}
\hypersetup{colorlinks=true,linkcolor=blue,citecolor=blue,urlcolor=blue}

\title{Steps used to build the analytic and skewer\\ Lyman-alpha $P_{1\mathrm{D}}$ pipelines}
\author{Helena}
\date{13 October 2026}

\begin{document}
\maketitle

\begin{abstract}
The task was a calculator. The inputs are the primordial parameters $n_s$ and $A_s$ and a cold-dark-matter spectrum $P_\mathrm{CDM}$. The output is the one-dimensional Lyman-alpha flux power $P_{1\mathrm{D}}(k_\parallel,z)$, with $k_\parallel$ in $\mathrm{s\,km}^{-1}$ and $P_{1\mathrm{D}}$ in $\mathrm{km\,s}^{-1}$. Two flux models share that matter spectrum. The analytic path is the Kaiser formula and a line-of-sight integral. The skewer path draws one-dimensional lines, applies the fluctuating Gunn--Peterson approximation, and measures the flux power. This note records every step, the equation used at that step, the assumptions, and the references. The cross-power $P_x$, the covariance, and the Fisher matrix are not in this calculation.
\end{abstract}

\section{What is delivered}
\label{sec:task}

One evaluation is at one redshift. The fiducial point used for the checks is
\begin{equation}
\begin{aligned}
\Omega_m &= 0.31, &
\Omega_b &= 0.049, &
h &= 0.67, \\
n_s &= 0.965, &
A_s &= 2.1\times 10^{-9}, &
z &= 3.
\end{aligned}
\end{equation}
These cosmological numbers are a rounded Planck-like choice, held fixed so that every check uses the same input. $P_\mathrm{CDM}$ can be built in three ways: the Eisenstein--Hu transfer function, a table read from a file (CLASS or CAMB), or the Eisenstein--Hu spectrum multiplied by a four-parameter bend.

Both pipelines then answer the same question: given that matter spectrum, what is the one-dimensional flux power. They do it at different points. The analytic path converts matter to flux before the line-of-sight integral. The skewer path integrates the matter spectrum first, builds lines, and only then converts density to flux.

\section{Shared step: the linear matter spectrum}
\label{sec:matter}

\subsection{Eisenstein--Hu transfer function}

The default three-dimensional spectrum is linear cold dark matter with the Eisenstein and Hu fitting formula \cite{Eisenstein:1997ik}. Wavenumbers passed into the transfer function are in $\mathrm{Mpc}^{-1}$, without $h$. The fit is normalised so that
\begin{equation}
T(k\to 0) = 1,
\end{equation}
and it includes the baryon acoustic oscillations. The internal fitting expressions for $T(k)$ are those of Eisenstein and Hu; they are not re-derived here.

\subsection{Primordial amplitude and growth}

The linear power in $\mathrm{Mpc}^{3}$ is
\begin{equation}
\label{eq:pcdm}
P(k,z)
=
\frac{8\pi^{2}}{25}
\frac{A_s}{\Omega_m^{2} H_0^{4}}
\left(\frac{k}{k_\mathrm{pivot}}\right)^{n_s-1}
k\,T(k)^{2}\,D(z)^{2}.
\end{equation}
The pivot is the standard primordial pivot,
\begin{equation}
k_\mathrm{pivot} = 0.05\,\mathrm{Mpc}^{-1}.
\end{equation}
$H_0$ is expressed in $\mathrm{Mpc}^{-1}$,
\begin{equation}
H_0 = \frac{100\,h}{299792.458}.
\end{equation}
Multiplication by $h^{3}$ converts Eq.~\eqref{eq:pcdm} to $(\mathrm{Mpc}/h)^{3}$, which is the unit used everywhere below. The wavenumber in that unit is $k_h = k/h$.

The growth factor $D(z)$ is the integral solution of the linear growth equation in a flat universe with constant dark energy. It is normalised so that $D(a)\to a$ deep in matter domination. It is not normalised to $D(0)=1$. The Hubble rate used for growth and for the later unit conversion is
\begin{equation}
\label{eq:hubble}
H(z) = 100\,h\sqrt{\Omega_m(1+z)^{3}+(1-\Omega_m)}
\quad[\mathrm{km\,s}^{-1}\,\mathrm{Mpc}^{-1}].
\end{equation}
Neutrino masses are not included. The universe is flat, with $\Omega_\Lambda = 1-\Omega_m$.

\subsection{A table instead of the fitting formula}

A whitespace-separated file with two columns may replace Eq.~\eqref{eq:pcdm}. Column one is $k$ in $h\,\mathrm{Mpc}^{-1}$. Column two is the linear $P(k)$ in $(\mathrm{Mpc}/h)^{3}$ at one fiducial redshift, already containing the fiducial $n_s$ and $A_s$. At least four strictly increasing positive rows are required, and $P>0$. Values are interpolated with a monotone cubic interpolant in $\ln k$ and $\ln P$. No extrapolation is allowed: a requested $k$ outside the file stops the calculation.

A later change of amplitude, tilt, or redshift rescales the table,
\begin{equation}
\label{eq:table}
P(k)
=
P_\mathrm{fid}(k)
\left(\frac{A_s}{A_{s,\mathrm{fid}}}\right)
\left(\frac{k}{k_p}\right)^{n_s-n_{s,\mathrm{fid}}}
\left(\frac{D(z)}{D(z_\mathrm{fid})}\right)^{2}.
\end{equation}
The pivot in $h\,\mathrm{Mpc}^{-1}$ is $k_p = 0.05/h$, so that it matches $k_\mathrm{pivot}=0.05\,\mathrm{Mpc}^{-1}$. If the file is already at the redshift of the call, the growth ratio is $1$.

CLASS often writes $k$ in $\mathrm{Mpc}^{-1}$ and $P$ in $\mathrm{Mpc}^{3}$. Those files have to be converted before reading,
\begin{equation}
k_h = k/h, \qquad P_h = P\,h^{3}.
\end{equation}
CAMB's matter-power output is already in $h$ units. The file must be the linear spectrum, not the nonlinear one.

\subsection{Optional four-parameter bend}

The task also asked for an $\alpha$--$\beta$--$\gamma$--$\delta$ shape. It is a smooth multiplier on the Eisenstein--Hu spectrum,
\begin{equation}
\label{eq:abgd}
P(k) = P_\mathrm{EH}(k)\,
\exp\!\left(\alpha L + \beta L^{2} + \gamma L^{3} + \delta L^{4}\right),
\qquad
L = \ln(k/k_p),
\end{equation}
with $k_p = 1\,h\,\mathrm{Mpc}^{-1}$. At that pivot the exponential equals $1$, so the amplitude remains the Eisenstein--Hu amplitude. Setting all four coefficients to zero returns $P_\mathrm{EH}$ exactly. These coefficients are shape parameters. They are not the IGM temperature--density index $\gamma$ of Sec.~\ref{sec:fgpa}. The bend does not add a sharp small-scale cutoff.

\section{Analytic pipeline}
\label{sec:kaiser}

\subsection{Kaiser flux power}

On large scales the forest is treated as a linearly biased tracer of the dark matter, with coherent peculiar velocities \cite{Kaiser:1987qv,McDonald:2001fe}. The three-dimensional flux power is
\begin{equation}
\label{eq:kaiser}
P_F(k,\mu)
=
b_F^{2}\,(1+\beta_F\mu^{2})^{2}\,P_\mathrm{CDM}(k).
\end{equation}
The direction cosine is
\begin{equation}
\mu = k_\parallel / k.
\end{equation}
The two flux parameters are fixed numbers,
\begin{equation}
b_F = -0.14, \qquad \beta_F = 1.4.
\end{equation}
$b_F$ is negative because a denser patch absorbs more light and transmits less. In Kaiser's derivation $\beta = f/b$, with $f$ the growth rate. This calculation does not evaluate that ratio. The pair $(-0.14,\,1.4)$ is a round stand-in for the large-scale forest bias. Early DESI data give $b_{\mathrm{Ly}\alpha}=-0.134\pm 0.009$ and $\beta_{\mathrm{Ly}\alpha}=1.41^{+0.12}_{-0.15}$ from the combined auto- and cross-correlation \cite{Gordon:2023dua}. The same parameters are held fixed at every redshift: the evolution $b(z)\propto (1+z)^{2.9}$ used in the eBOSS and DESI analyses \cite{eBOSS:2020tmo} is not applied. McDonald's small-scale corrections for nonlinear growth, gas pressure, and nonlinear velocities \cite{McDonald:2001fe} are not applied either. Equation~\eqref{eq:kaiser} is linear Kaiser only.

This path never builds a density and never evaluates an optical depth. Equation~\eqref{eq:kaiser} is the small-fluctuation limit of the optical-depth model in Sec.~\ref{sec:fgpa}, not a fit of $b_F$ to that model. With the skewer normalisation below, the linear response would be closer to $b_F\sim -A\beta\approx -0.8$. The value $-0.14$ is the measured large-scale bias, after the smoothing that the skewers leave out.

\subsection{Line-of-sight integral}

A forest spectrum sees only the line-of-sight direction. Perpendicular modes are integrated out,
\begin{equation}
\label{eq:p1d-kaiser}
P_{1\mathrm{D}}(k_\parallel)
=
\int_{k_\parallel}^{k_\mathrm{max}}
\frac{k\,\mathrm{d}k}{2\pi}\,
P_F\!\left(k,\mu=\frac{k_\parallel}{k}\right).
\end{equation}
The upper limit is a numerical cutoff,
\begin{equation}
k_\mathrm{max} = 30\,h\,\mathrm{Mpc}^{-1},
\end{equation}
sampled with 800 logarithmically spaced points and a trapezoidal rule. The integral starts at $k_\parallel$ because a three-dimensional mode cannot contribute to a line-of-sight wavenumber larger than its own $k$. The factor $k\,\mathrm{d}k/(2\pi)$ is the ring of perpendicular modes at fixed $k_\parallel$.

$k_\mathrm{max}$ is not a measured forest scale. The one-dimensional forest power used by BOSS, eBOSS, and DESI, and in the check plots, spans
\begin{equation}
k_\parallel \sim 10^{-3}\ \mathrm{to}\ 0.02\ \mathrm{s\,km}^{-1}.
\end{equation}
At $z=3$ in the fiducial cosmology, $H(z)/[(1+z)h]\approx 113\,\mathrm{km\,s}^{-1}$ per $\mathrm{Mpc}/h$, so that range is about $0.11$ to $2.3\,h\,\mathrm{Mpc}^{-1}$, or wavelengths of about $55$ down to $3\,\mathrm{Mpc}/h$. The integral extends past those points because small-scale power still contributes to a lower $k_\parallel$. Raising $k_\mathrm{max}$ until $P_{1\mathrm{D}}$ stops changing is the practical test of the cutoff. Gas pressure removes power around $10$--$20\,h\,\mathrm{Mpc}^{-1}$ \cite{McDonald:2001fe}; integrating the unsmoothed Eq.~\eqref{eq:kaiser} out to $30\,h\,\mathrm{Mpc}^{-1}$ keeps power that the gas does not have.

\subsection{Forest units}

The conversion between comoving $h\,\mathrm{Mpc}^{-1}$ and the spectrograph velocity unit is
\begin{equation}
\label{eq:units}
\frac{\mathrm{d}v}{\mathrm{d}\chi}
=
\frac{H(z)}{(1+z)\,h}
\quad[\mathrm{km\,s}^{-1}\ \mathrm{per}\ \mathrm{Mpc}/h],
\end{equation}
with $H(z)$ from Eq.~\eqref{eq:hubble}. Then
\begin{equation}
k_h = k_{\mathrm{s/km}}\times \frac{\mathrm{d}v}{\mathrm{d}\chi},
\qquad
P_{1\mathrm{D}}[\mathrm{km\,s}^{-1}]
=
P_{1\mathrm{D}}[\mathrm{Mpc}/h]\times \frac{\mathrm{d}v}{\mathrm{d}\chi}.
\end{equation}
The analytic output is this $P_{1\mathrm{D}}$ in $\mathrm{km\,s}^{-1}$.

\subsection{Where the Kaiser formula stops holding}

The formula assumes a small density contrast, so that the flux contrast is proportional to $\delta$, and a coherent infall, which is the regime Kaiser derived. It fails in four places that matter for the forest.
\begin{itemize}
\item The map $F=e^{-\tau}$ is linear only for a small $\delta$. Saturated lines stop responding once the pixel is already dark.
\item Random thermal and turbulent motions suppress power at $\mu\sim 1$ on small scales. That suppression is in McDonald's correction \cite{McDonald:2001fe} and is absent here.
\item Pressure and thermal broadening remove small-scale power. The bare Kaiser spectrum does not.
\item One constant $b_F$ cannot describe strong overdensities or gas near galaxies, where the optical-depth rule itself fails \cite{Kooistra:2022hzw}.
\end{itemize}
The shape of the response to $n_s$ and $A_s$ on the largest forest scales is what this path is for. The absolute $P_{1\mathrm{D}}$ near $0.02\,\mathrm{s\,km}^{-1}$ uses the formula outside that regime.

\section{Skewer pipeline}
\label{sec:skewers}

The skewer path uses the same $P_\mathrm{CDM}$ and then builds explicit lines of sight. It does not build a three-dimensional cube.

\subsection{From the three-dimensional spectrum to one line}

A Gaussian cube with power $P_\mathrm{CDM}(k)$ would be filled in Fourier space and transformed back, and a skewer would be one row of that cube. For an isotropic spectrum the power along one row is fixed before the cube exists,
\begin{equation}
\label{eq:p1d-matter}
P_{1\mathrm{D},\mathrm{matter}}(k_\parallel)
=
\frac{1}{2\pi}
\int_{k_\parallel}^{q_\mathrm{max}}
q\,P_\mathrm{CDM}(q)\,\mathrm{d}q.
\end{equation}
There is no $b_F$ and no $\mu$ in this integrand. The upper limit is the last $k$ in the matter table, not the $30\,h\,\mathrm{Mpc}^{-1}$ of Eq.~\eqref{eq:p1d-kaiser}. The lines are then drawn in one dimension so that they carry Eq.~\eqref{eq:p1d-matter}.

Statistically, one such line matches one row cut from an isotropic Gaussian cube. Neighbouring rows in a cube pass through the same clumps. These lines are drawn independently, so they are uncorrelated, and there is no transverse velocity field.

\subsection{Drawing one line}

A line is periodic, of length $L$, with $N$ pixels. The pixel size and the highest resolved wavenumber are
\begin{equation}
\mathrm{d}x = L/N,
\qquad
k_\mathrm{Nyquist} = \pi\,N/L.
\end{equation}
The check in the summary figure uses $L=200\,\mathrm{Mpc}/h$ and $N=512$, so $\mathrm{d}x=0.391\,\mathrm{Mpc}/h$ and $k_\mathrm{Nyquist}\approx 8\,h\,\mathrm{Mpc}^{-1}$. At $z=3$ that pixel is about $44\,\mathrm{km\,s}^{-1}$ wide. The longest wave on the line is $2\pi/L$. The matter table must cover from that fundamental mode through the Nyquist mode.

The line is not filled pixel by pixel. Each Fourier mode with $k>0$ is a complex Gaussian number whose variance is set by $P_{1\mathrm{D},\mathrm{matter}}(k)$. The constant mode is set to zero, so the line fluctuates about zero. An inverse Fourier transform gives the Gaussian contrast $\delta_G(x)$. The random seed only makes the draw repeatable. The plotted check uses 256 lines and seed 7. The noise on the measured flux power falls as $1/\sqrt{N_\mathrm{lines}}$.

An optional Gaussian smooth of $\delta_G$, with width given in pixels, stands in for pressure smoothing. Its default width is zero.

\subsection{Lognormal density and optical depth}
\label{sec:fgpa}

The Gaussian contrast is mapped to a positive density with mean near one \cite{Bi:1996fh,Coles:1991if},
\begin{equation}
\label{eq:lognormal}
\Delta = \exp\!\left(\delta_G - \tfrac{1}{2}\mathrm{Var}(\delta_G)\right).
\end{equation}
The optical depth follows the fluctuating Gunn--Peterson approximation \cite{Gunn:1965hd,Bi:1996fh,Hui:1997dp}. Photoionization equilibrium and the temperature--density relation $T\propto\Delta^{\gamma-1}$, together with a recombination rate $\propto T^{-0.7}$, give
\begin{equation}
\label{eq:beta}
\beta = 2 - 0.7\,(\gamma - 1),
\qquad
\tau = A\,\Delta^{\beta},
\qquad
F = e^{-\tau}.
\end{equation}
The fiducial temperature--density index is $\gamma=1.5$, so $\beta=1.65$. This $\beta$ is the IGM exponent. It is not the Kaiser parameter $\beta_F$. Kooistra, Lee, and Horowitz use the same exponent \cite{Kooistra:2022hzw}. Arya et al.\ use the same lognormal-to-flux sequence, with a Jeans cutoff and an amplitude set by temperature and the photoionization rate \cite{Arya:2022uke}.

$A$ is not taken from that physical expression. It is solved so that the mean transmitted flux, averaged over every pixel of every line, equals a chosen target,
\begin{equation}
\label{eq:meanflux}
\left\langle e^{-A\,\Delta^{\beta}} \right\rangle = \bar F.
\end{equation}
The target used in the checks is $\bar F = 0.80$. The root is unique because the mean flux falls as $A$ increases. For the plotted lines the solution is $A=0.485$. The target is an input, not a function of redshift. At $z=3$ the observed mean transmission is somewhat lower, closer to $0.7$.

Peculiar velocities are applied only if an external pixel displacement is supplied. The displacement remaps $\tau$ with periodic cloud-in-cell assignment. The checks do not supply one, so those lines are real space. An optional Gaussian smooth of $\tau$, also off by default, stands in for thermal broadening. It is not a Voigt profile, and it does not vary from pixel to pixel.

\subsection{Measured flux power}

The flux contrast on the lines is
\begin{equation}
\delta_F = F/\langle F\rangle - 1.
\end{equation}
Its one-dimensional power is the average, over the lines, of the squared Fourier amplitude,
\begin{equation}
\label{eq:measure}
P_{1\mathrm{D}}(k)
=
\frac{(\mathrm{d}x)^{2}}{L}
\left\langle \left|\tilde\delta_F(k)\right|^{2} \right\rangle,
\end{equation}
in $\mathrm{Mpc}/h$. The same factor Eq.~\eqref{eq:units} converts $k$ to $\mathrm{s\,km}^{-1}$ and $P_{1\mathrm{D}}$ to $\mathrm{km\,s}^{-1}$. This spectrum is a measurement, not another integral of a three-dimensional flux power.

\subsection{Density variance and the papers' caveats}

With pressure smoothing left off, the lognormal density variance on these lines is about $19$. That is far outside the small-fluctuation assumption behind Eq.~\eqref{eq:kaiser}. Arya et al.\ suppress the small-scale matter power with a Jeans factor before the lognormal map \cite{Arya:2022uke}. The switch for that smoothing exists and is off. Kooistra, Lee, and Horowitz show that $\tau\propto\Delta^{\beta}$ already fails in strong overdensities and where galactic feedback matters \cite{Kooistra:2022hzw}. Neither effect is modelled.

\section{Why the two flux spectra differ}
\label{sec:compare}

On the summary figure the analytic curve and the skewer curve are both $P_{1\mathrm{D}}$ in $\mathrm{km\,s}^{-1}$ at $z=3$. They are not expected to lie on top of each other.
\begin{itemize}
\item The analytic height is $b_F^{2}$ with $b_F=-0.14$, times the velocity factor. The skewer height is fixed by $\bar F=0.80$ through Eq.~\eqref{eq:meanflux}.
\item The analytic spectrum includes peculiar velocities. The plotted skewers are real space.
\item The analytic spectrum is linear. The skewers use the exponential $e^{-\tau}$ on a field with density variance about $19$.
\item The skewer curve is an average of 256 finite lines, so it is noisy. The analytic curve is not random.
\end{itemize}
Dividing each curve by its own value at $0.005\,\mathrm{s\,km}^{-1}$ removes the normalisation. The shapes then follow each other. For a small, smoothed, real-space field the two would agree in amplitude as well, with $b_F=-A\beta$ and $\beta_F=0$. Matching the factor $(1+\beta_F\mu^{2})^{2}$ requires velocities that know about modes perpendicular to the line of sight. Independent one-dimensional lines do not have those modes.

A one percent increase in $A_s$ raises the analytic $P_{1\mathrm{D}}$ by one percent at every $k_\parallel$. A change in $n_s$ tilts it. That flat amplitude response is the property a later Fisher forecast needs, which is why derivatives would be taken on the analytic path. The skewer power depends on the random seed.

\section{Assumptions, collected}
\label{sec:assumptions}

\begin{enumerate}
\item Flat $\Lambda$CDM, no massive neutrinos, linear matter power only.
\item Eisenstein--Hu $T(k)$ with $T(0)=1$, including baryon oscillations \cite{Eisenstein:1997ik}. It is a fit, not a Boltzmann solution. A CLASS or CAMB table can replace it through Eq.~\eqref{eq:table}.
\item Growth normalised by $D(a)\to a$ in matter domination, not by $D(0)=1$.
\item One redshift per evaluation. $b_F$, $\beta_F$, $\gamma$, and $\bar F$ do not evolve inside a call.
\item Analytic flux power is linear Kaiser, Eq.~\eqref{eq:kaiser}, with fixed $b_F=-0.14$ and $\beta_F=1.4$ \cite{Kaiser:1987qv,McDonald:2001fe,Gordon:2023dua}. No McDonald small-scale correction.
\item The analytic integral stops at $30\,h\,\mathrm{Mpc}^{-1}$ by choice.
\item Skewers are independent periodic lines, not rows of one cube. No transverse correlations.
\item Real space, unless an external displacement is passed in. Pressure and thermal smoothing default to zero.
\item Lognormal density, Eq.~\eqref{eq:lognormal} \cite{Bi:1996fh}. Optical depth Eq.~\eqref{eq:beta} with $\gamma=1.5$ \cite{Hui:1997dp,Kooistra:2022hzw}. $A$ is fit to $\bar F=0.80$, not computed from $T_0$ and the photoionization rate \cite{Arya:2022uke}.
\item The $\alpha\beta\gamma\delta$ coefficients in Eq.~\eqref{eq:abgd}, if used, are a smooth shape and not the IGM $\gamma$.
\item Tables are not extrapolated.
\end{enumerate}

\section{What this calculation does not contain}

The cross-power of the forest with quasars or galaxies, the covariance, survey area, number densities, resolution, noise, and the Fisher matrix are the next stage. The covariance of the one-dimensional flux power across redshift bins, including the effect of the longest waves along a sightline, is studied with this same lognormal model by Arya, Paranjape and Choudhury \cite{Arya:2023ipe}. LaCE, LyaCoLoRe, cup1d, P1D\_Mocks, and qsotools are not part of the calculator. LaCE emulates hydrodynamic spectra. LyaCoLoRe builds a three-dimensional mock and is the relevant comparison when correlated sightlines are required. P1D\_Mocks and qsotools start from a $P_{1\mathrm{D}}$ that already exists.

\bibliographystyle{unsrtnat}
\bibliography{references}

\end{document}
